\subsubsection{Resumable Gather and Scatter}

(:ref:VECTOR.instructions.VGATHER:) and (:ref:VECTOR.instructions.VSCATTER:)
are full-vector memory instructions with an explicit completion predicate.
Their canonical operand orders are illustrated below; $T$ denotes the selected
\texttt{B/W/L/Q} element width.

\begin{verbatim}
VGATHER.T  Pg, Ps, [address], Vd
VSCATTER.T Pg, Ps, Vs, [address]
\end{verbatim}

\paragraph{Typed completion image.}
For element width $E$ bytes, let $N=VLEN/(8E)$.  Define $G=typed_E(Pg)$ by
copying \texttt{Pg[iE]} to bit $iE$ for every $0\leq i<N$ and clearing every
other predicate bit.  At each entry or re-entry, the instruction performs the
architectural normalization

\[
  Ps \leftarrow typed_E(Ps) \mathbin{\&} G,
  \qquad pending = G \mathbin{\&} \mathop{\sim}Ps.
\]

\texttt{Pg} and \texttt{Ps} shall be distinct.  The normalization is committed
before a lane memory transaction is issued and remains visible if that
transaction faults.  Consequently inactive completion bits and all
nonsignificant bits are zero.  If \texttt{pending} is empty, the instruction
retires without a memory transaction and its normal postcondition is
\texttt{Ps == G}.

Otherwise the implementation selects any pending lane $i$.  VGATHER performs
one $E$-byte load and, on success, writes only \texttt{Vd[i]} before setting
\texttt{Ps[iE]}.  VSCATTER performs one $E$-byte store of \texttt{Vs[i]} and,
on success, sets \texttt{Ps[iE]}.  Every successful lane therefore has one
completion bit that journals its externally visible work.  VGATHER requires
its vector address register to differ from \texttt{Vd}; VSCATTER permits its
vector address register and \texttt{Vs} to be the same register.

\paragraph{Lane-completion boundary.}
After a successful lane, destination or memory state and the corresponding
completion bit commit together.  If another lane is pending, the PC remains at
the gather or scatter instruction and an asynchronous event may be admitted.
This progress boundary is not architectural instruction retirement, does not
increment an architectural retired-instruction count, and does not itself
generate a debug-trace event.  The final empty-pending transition advances the
PC and forms the (:term:base.terms.ordinary:) retirement and trace boundary.

An event handler may emulate or waive a faulting lane by setting its completion
bit.  For VGATHER, emulation shall first write the lane's destination value;
for VSCATTER, it shall first perform or explicitly waive the store.  Returning
to the saved PC is sufficient to continue: no hidden cursor or implementation
token is part of (:term:base.terms.architectural_state:).

\paragraph{Address forms.}
The lane number $i$ is implicit and does not appear in assembly syntax.  Before
default-data-segment translation, the address forms are:

\begin{center}
\begin{tabular}{lll}
\toprule
Canonical address & Constraint & Lane offset \\
\midrule
\texttt{[Rb + Rs]} & any $T$ & $Rb+i\mathbin{\cdot}signed64(Rs)$ \\
\texttt{[Va]} & \texttt{L} & $zero\_extend_{64}(Va.L[i])$ \\
\texttt{[Va]} & \texttt{Q} & $Va.Q[i]$ \\
\texttt{[Rb + Va]} & any $T$ & $Rb+sign\_extend_{64}(Va.T[i])$ \\
\texttt{[Rb + Va * scale]} & any $T$ & $Rb+zero\_extend_{64}(Va.T[i])E$ \\
\texttt{[Rb + Va * scale + disp]} & any $T$ & preceding indexed address plus \texttt{disp} \\
\bottomrule
\end{tabular}
\end{center}

In the scalar-stride form, \texttt{Rs} supplies the signed lane stride even
though the canonical spelling is the familiar-looking \texttt{[Rb + Rs]}.
In vector-index forms, the assembler token \texttt{scale} is fixed by $T$ to
$E$; it is not an additional operand.  The \texttt{disp8s}, \texttt{disp16s},
and \texttt{disp32s} (:term:base.terms.payload|plural:) are sign-extended, while \texttt{disp64} contributes
its exact 64-bit bit pattern.  All address arithmetic is modulo $2^{64}$.  The
absolute forms are still offsets through the default data segment and do not
bypass segmentation.

\paragraph{Ordering, overlap, and failures.}
Lane selection and memory-access order are not architecturally specified.
Each successful lane access is nevertheless one (:term:base.terms.ordinary:), naturally sized
$E$-byte normal-memory access and obeys the baseline memory-ordering rules.
If two VSCATTER lanes target overlapping physical bytes, which lane supplies
the final value of each overlapping byte is unspecified.  This does not merge,
drop, or widen either access: each successful lane still performs its own
$E$-byte store and sets its own completion bit.

A load fault or a scatter failure reported before the store point of no return
leaves the selected lane pending and raises the (:term:base.terms.ordinary:) precise synchronous
fault.  An (:term:base.terms.ordinary:) memory bus failure in this state is retry-safe.  If a
scatter store has become irrevocable, the implementation sets the lane's
completion bit and shall not report a restartable synchronous failure for that
lane.  Any later report is an imprecise, non-retry-safe
(:ref:base.events.EXCEPTION.MACHINE_CHECK:) at a subsequent lane or retirement
boundary.  These instructions are not permitted for MMIO, including when only
one lane is active.  Both preserve FLAGS.
